\begin{abstract}
% As the workloads demanded of LLMs grow more complex, the length of outputs generated by LLMs has increased, and the bottleneck in LLM inference is progressively shifting towards the decode stage. Previous studies have primarily employed two main methods to enhance the throughput of the decode stage on xPUs. One is to increase the arithmetic intensity of the decode stage through batching, and the other is to reduce the data loaded from memory by utilizing lossy compression. Especially in edge LLM inference scenarios, where batching cannot be effectively utilized, lossy compression becomes key to decode stage acceleration.

% However, with the recent emergence of reasoning LLMs, the significant accuracy degradation introduced by lossy compression methods has become a major problem. This accuracy degradation, even as sLLMs are now incorporating reasoning capabilities, prevents us from fully leveraging the potential of modern LLMs.

% Due to these challenges, LLM acceleration without accuracy degradation becomes increasingly important, and speculative decoding can be a key solution to meet this demand. In contrast, existing speculative decoding models suffer from limitations that make them unsuitable for edge device inference, such as high computational resource consumption, large memory capacity requirements, and low performance improvement in small batchs.

% To address these issues, we propose Cassandra. Cassandra is a cost-effective and easy-to-apply algorithm-hardware co-designed speculative decoding solution for anyone seeking to perform and accelerate LLM inference on consumer-grade GPUs and NPUs. Cassandra is completely training-free, requires no additional memory capacity, and demonstrates superior performance across various tasks, including reasoning tasks and long contexts. Cassandra can be utilized with only minor hardware extensions to existing xPUs, and the overhead is negligible. Our experimental results show that this approach yields up to 2.37$\times$ performance improvement on various benchmarks without any accuracy loss.

% However, despite the introduction of various speculative decoding techniques, a method that can provide sufficient performance at low batch sizes without necessitating additional training remains elusive. This presents a considerable challenge for achieving lossless acceleration in low-batch LLM inference on consumer-grade devices.
% To address this limitation, we propose Cassandra, an algorithm-hardware co-designed self-speculative decoding method tailored explicitly for low-batch scenarios. The core idea of Cassandra is to create a high-performance training-free draft model by leveraging fine-grained data selection criteria. Utilizing optimized pruning and mantissa truncation, Cassandra isolates the most salient values within the weights and the Key-Value (KV) cache from the remaining data. Subsequently, only these most salient values are initially loaded to rapidly generate a few tokens, after which the full-precision values are loaded to verify this result in parallel. This approach yields a substantial performance enhancement compared to existing self-speculative decoding methods, which typically rely on techniques such as layer-skipping or structured KV cache compression. Furthermore, to mitigate the potential bottleneck of format conversion between the Cassandra format and standard floating-point representation, we introduce a lightweight encoder and decoder hardware designed for integration and use on commercial GPUs and NPUs.

Speculative decoding has emerged as a promising lossless approach to accelerating Large Language Models (LLMs). Recently, as the overhead of the decode stage has increased in reasoning LLMs, while algorithmic approximations cause significant accuracy degradation, lossless acceleration through speculative decoding has become a key component for efficient LLM serving systems. However, despite numerous advances, a speculative decoding method that can provide sufficient performance at low batch sizes without requiring additional training remains elusive. This presents a considerable challenge for achieving lossless acceleration in low-batch LLM inference on consumer-grade devices.

To address this limitation, we propose Cassandra, an algorithm–hardware co-designed self-speculative decoding framework tailored explicitly for low-batch scenarios. The core idea of Cassandra is to construct a high-performance, training-free draft model through fine-grained data selection. By leveraging optimized pruning and mantissa truncation, Cassandra identifies and isolates the most salient values within both model weights and the Key–Value (KV) cache. These selected values are first loaded to rapidly generate candidate tokens, followed by parallel verification using full-precision data. This design achieves substantial performance gains over prior self-speculative decoding methods, which primarily rely on layer skipping or structured KV cache compression. In addition, to mitigate the overhead of format conversion between the Cassandra representation and standard floating-point formats, we introduce a lightweight encoder–decoder hardware module designed for seamless integration with commercial GPUs and NPUs.

Our evaluation demonstrates that Cassandra achieves a performance improvement of up to 2.41$\times$ than the BFloat16 baseline without additional training. Moreover, when running Llama3-8B on an RTX 4090, Cassandra is capable of generating 1.81$\times$ more tokens under a fixed memory budget compared to Eagle-3, a state-of-the-art speculative decoding method.

\end{abstract}